% RESEARCH_DRAFT_SHA256: e974d1019e3676511c9f6cc2282893998aee50e1f3ad589e69083504d8b9c640
% Layout and source-reference wrappers expanded for the web renderer.
\documentclass[11pt]{article}
\usepackage[margin=25mm]{geometry}
\usepackage{fontspec,unicode-math}
\setmainfont{Latin Modern Roman}
\setmathfont{Latin Modern Math}
\usepackage{amsmath,amssymb,mathtools,xurl,hyperref}
\hypersetup{colorlinks=true,urlcolor=blue,linkcolor=blue}
\setcounter{secnumdepth}{0}
\setlength{\parindent}{0pt}
\setlength{\parskip}{5pt}
\setlength{\emergencystretch}{3em}
\newcommand{\sref}[2]{\hyperlink{source:#1:#2}{[S#2]}}
\newcommand{\cataloguescope}[1]{\paragraph{Source record.}#1}
\newcommand{\floor}[1]{\left\lfloor#1\right\rfloor}
\begin{document}
\section*{Research-attempt provenance}
This standalone research attempt retains the complete original definition body below,
including its historical source and preparation notes. Those notes describe the
earlier catalogue preparation; the new attempt and its verification are identified
separately after that body. No original ID, catalogue status, or source statement is changed.

\section{An intrinsic-dimension certificate for projected Hamilton--Jacobi equations}\label{statement}

\subsection{Definitions and mathematical statement}

\paragraph{Shared notation.}Standard number systems are $\mathbb N,\mathbb Z,\mathbb Q,\mathbb R,\mathbb C$. Specialized notation and index ranges follow the cited source and the definitions below.


A Hamilton--Jacobi equation has the form $\partial_tV+H(x,\nabla V)=0$ for a specified Hamiltonian $H$. Level-set methods evolve a scalar function whose level sets represent geometric fronts or controlled reachable sets. High dimension increases the number of grid points needed by tensor-product discretizations.
\paragraph{Record scope and qualifications.}
The catalogue title is retained as metadata. Named certificates, contracts, bounds or frameworks in that title are not asserted as results of the original source extract.
This brief source record names an area of interest only. It states no projection, intrinsic-dimension certificate, Hamiltonian, or numerical guarantee. The supplied catalogue marks this record open; its present research status was not independently checked.
The source text below is rendered literally where necessary; damaged formulas and composite extracts are not silently repaired into a new mathematical claim.
\paragraph{Statement recorded in the supplied source.}
A.9 Osher, Stanley

I am interested in solving high dimensional Hamilton-Jacobi equations numerically, especially those arising in control theory and involving the level set method. \href{https://aimath.org/WWN/optimalcontrol/optimalcontrol.pdf}{[S2]}
\paragraph{Source status.}
The supplied catalogue labels this source record open; that metadata was checked against the supplied snapshot on 2026-10-11. The label is preserved for provenance and is not an independent verification of present open status or mathematical correctness.
\subsection{Short English statement}
The participant expresses interest in numerical Hamilton--Jacobi equations and level-set methods.
\subsection{Sources}
\begin{itemize}
\item[S1] Ulam AI and the UnsolvedMath Contributors, UnsolvedMath: A Curated Collection of Open Mathematics Problems, catalogue entry. Catalogue curation and metadata: CC BY 4.0. The source extract is retained; definitions, English summaries and qualifications are editorial additions.
\url{https://huggingface.co/datasets/ulamai/UnsolvedMath}.
\item[S2] AIM Workshop Problem Lists, source linked by the catalogue.
\url{https://aimath.org/WWN/optimalcontrol/optimalcontrol.pdf}.
\end{itemize}
\paragraph{Preparation and review.}
The original source and UnsolvedMath contributors retain attribution for the source extract. Codex (GPT-6; exact serving revision not exposed) assisted with definitions, English summaries, scope notes and formatting. No new human mathematical authorship or independently verified solution is asserted. Catalogue identity checked against the complete supplied registry and current repository. Distinct upstream records, including repeated questions and broad or damaged extracts, are retained for complete coverage of the source CS category. A repeated record is not a new independent conjecture. Complete mathematical-equivalence and current-literature review has not been performed. Preparation and final review follow contribution prompts C and D, with source gaps disclosed. Structural and rendering checks do not establish a new theorem or open status.
\label{end:statement}

\section{Literature and context check}
The participant contribution in the primary AIM optimal-control workshop
 PDF was consulted. Its broad scope is retained above. The following explicitly
 chosen finite or smooth model tests one route suggested by the agenda; it is
 not a reconstruction of a missing theorem. The elementary arguments are not
 claimed novel, and no current-resolution audit of the whole field is asserted.
\section{Research attempt: A projected Hamiltonian must retain the projection metric}
For a smooth auxiliary HJ equation $u_t+H(\nabla_xu)=0$, try the ansatz
$u(t,x)=v(t,Px)$ for a fixed linear map $P:\mathbb R^d\to\mathbb R^k$.
The chain rule gives the exact reduced equation
$v_t+H(P^T\nabla_yv)=0$. For $H(p)=\|p\|^2/2$, its reduced Hamiltonian is
$q^TPP^Tq/2$, not generally $\|q\|^2/2$. Thus even a correct ridge ansatz
does not permit ignoring the geometry of the projection.

Take $d=2,k=1$, $P(x)=2x_1$ and initial data $u(0,x)=2x_1$.
The exact smooth solution is $u(t,x)=2x_1-2t$: its time derivative is $-2$
and gradient norm squared is four. Reusing the unit-metric reduced Hamiltonian
would instead give $v(t,y)=y-t/2$, and hence $u(t,x)=2x_1-t/2$, whose HJ
residual is $3/2$. This is a fully computed failure of the naive dimensional
reduction despite exactly ridge-shaped initial data.

For a state-dependent Hamiltonian $H(x,p)$, closure requires
$H(x,P^Tq)$ to depend on $x$ only through $Px$ for every relevant $q$.
This is sufficient for the displayed smooth substitution. It is also necessary
for a closed rule that works for every affine ridge test: the affine test
$v(y)=q\cdot y$ has the same reduced state and gradient on an entire fiber,
so the asserted reduced Hamiltonian must take one value there. For instance
$H(x,p)=\sin(x_2)p_1$ with $P(x)=x_1$ fails this test at $q=1$.

These chain-rule facts do not prove that arbitrary data become low-dimensional,
that viscosity test functions transfer under every map, or that a numerical
scheme avoids exponential complexity. Those are separate claims absent from
the very short source agenda.
\section{Verification and adversarial audit}
The scaled linear solution and residual were checked exactly. The metric
is unchanged when $PP^T=I$, explaining the special orthonormal-row case.
The statement is explicitly about smooth auxiliary solutions and test functions,
so no unproved viscosity equivalence is hidden in the substitution.
\paragraph{Audit corrections.} Projection scaling was retained, and a classical chain-rule identity was not promoted to a general high-dimensional numerical theorem.
\paragraph{External computation.} None; symbolic arguments and boundary checks only.
\section{Scope and remaining gaps}
The source supplies no Hamiltonian, initial data or error criterion. The
first gap is establishing an appropriate low-dimensional invariant structure
for actual data and transferring the relevant viscosity and numerical guarantees.
\paragraph{Integration review.} Codex (GPT-6; exact serving revision not exposed) reviewed the written arguments and expanded layout and source-reference wrappers. Repository metadata and the production web renderer were checked using Python 3.13 and JavaScript with MathJax 3.2.2. This model-assisted review is not independent human proof verification; no PDF compilation or formal proof checking is asserted.

\section{Completion Estimate}
\noindent\textbf{COMPLETION: 0\%}
The calculations below the preserved source statement investigate a concrete route,
but establish no substantiated new progress toward the entire original target.
Known special cases and failed sufficient conditions are not credited as a solution.
\section{References}
The original source references and catalogue attribution are preserved in the
Sources section above. The following URLs were actually consulted for this attempt;
retention of an additional original reference does not claim fresh consultation of it.
\begin{itemize}
\item \url{https://aimath.org/WWN/optimalcontrol/optimalcontrol.pdf} (consulted 2026-10-11). 
\end{itemize}
\section{Model and contributor disclosure}
OpenAI gpt-6-astra, requested reasoning configuration ultra, generated this research
attempt on 2026-10-11. The exact serving revision is not exposed. Codex GPT-6
performed orchestration, repository-format review, and a model-assisted mathematical
audit of the written arguments; its exact serving revision
is also not exposed. The mathematical route, boundary checks and final adversarial
self-audit followed contribution prompts A then B in
\texttt{docs/contribution\_prompts.md}, after reading \texttt{CONTRIBUTING.md}.
Python 3.14.2 assembled the standalone documents; any mathematical computation is
specified in Verification. No theorem prover was used. No human mathematical
authorship, substantial human proof correction, or independent human proof checking
is asserted. Model-assisted review, source retrieval, and structural checks do not
establish correctness or novelty. Full private reasoning transcripts are unavailable
as reproducibility artifacts; the complete arguments, inputs and outcomes actually
used are stated here or in the companion batch audit.
\section{Final status}
Unresolved as a new research contribution. This is a concrete route investigation,
not a claim of new complete resolution or a current-status certification. A known
answer reproduced for context does not become a new contribution. The catalogue's
source-status label is preserved; this attempt status does not assert that a
source-labelled solved question is open.
\end{document}
